Observemos, por otra parte, que $\operatorname{Spec}(A)$ consta del
punto cerrado $\mathfrak m$ y del punto genérico $(0)$, cuyos cuerpos
residuales son $k$ y $K$, respectivamente; como, en el caso en que nos hemos situado,
el cuerpo de fracciones de $B$ y su cuerpo residual son respectivamente isomorfos a
$K$ y $k$, las fibras del morfismo
$u:\operatorname{Spec}(B'_{\mathfrak p'})\to\operatorname{Spec}(B)$
sobre el punto cerrado y el punto genérico de $\operatorname{Spec}(B)$
son respectivamente \emph{isomorfas} a las fibras del morfismo
$\operatorname{Spec}(A'_{\mathfrak p'})\to\operatorname{Spec}(A)$
sobre el punto cerrado y el punto genérico de $\operatorname{Spec}(A)$
(\textbf{I}, 3.6.4), y por hipótesis son \emph{geométricamente normales}. Como además
$u$ es plano, se concluye que es \emph{normal} (6.8.1). Puesto que $B$ y
$B'_{\mathfrak p'}$ son noetherianos y $B$ es normal, se deduce de
(6.5.4) que $B'_{\mathfrak p'}$ es \emph{normal}.

